\subsection{Паттерны по типам запросов и архитектурные наблюдения}

Объединение результатов двух серий позволяет выделить устойчивые закономерности
(таблица~\ref{tab:patterns}). В потоковом сценарии сравнение ведётся по медиане как
менее зашумлённой метрике, отражающей типичную задержку, к которой и чувствителен
потребитель витрины данных.

{\small
\begin{xltabular}{\textwidth}{|>{\raggedright\arraybackslash}p{0.26\textwidth}|>{\raggedright\arraybackslash}p{0.18\textwidth}|>{\raggedright\arraybackslash}p{0.18\textwidth}|X|}
    \caption{Сравнение форматов по типам запросов и операций (NRT-сценарий, медиана)}\label{tab:patterns} \\ \hline
    \textbf{Тип операции} & \textbf{Vortex против Parquet} & \textbf{Vortex против ORC} & \textbf{Объяснение} \\ \hline
    \endfirsthead
    \caption*{Продолжение таблицы~\ref{tab:patterns}} \\ \hline
    \textbf{Тип операции} & \textbf{Vortex против Parquet} & \textbf{Vortex против ORC} & \textbf{Объяснение} \\ \hline
    \endhead
    \hline
    \endfoot
    \hline
    \endlastfoot
    Запись (INIT) & сопоставимы & сопоставимы & оба буферизуют строки и сбрасывают по чекпойнту, различия в пределах разброса \\ \hline
    Точечный доступ по ключу & сопоставимы & сопоставимы & фильтры Блума и статистики работают у всех; на маленьких выборках различия в пределах разброса \\ \hline
    Полное сканирование с проекцией & сопоставимы & Vortex чуть быстрее в FINAL ($\approx 1\,\%$) & column pruning + плотные кодеки работают на упорядоченных данных; на свежих L0-файлах эффект слабее \\ \hline
    Диапазонный фильтр по строке-префиксу & Vortex заметно быстрее (до $3{\times}$ в DURING) & Vortex стабильно быстрее ($\approx 4\text{--}25\,\%$) & FSST-кодирование строк с общими префиксами + нативное проталкивание предиката + пропуск файлов по полным статистикам \\ \hline
    Подсчёт с селективным предикатом & Vortex быстрее ($\approx 21\,\%$ в FINAL) & Vortex стабильно быстрее ($\approx 4\text{--}21\,\%$) & предикат отсекает данные до пересечения границы JNI; выигрыш от проталкивания преобладает над накладными расходами \\ \hline
    Многошаговая аналитика (соединения с фильтром, последовательности событий) & Vortex стабильно быстрее ($\approx 33\,\%$) & Vortex стабильно быстрее ($\approx 20\text{--}33\,\%$) & многократные проходы и предикаты, выигрыш от column pruning и проталкивания умножается \\ \hline
    Простая группировка / многомерная агрегация без предиката & Parquet иногда быстрее (единицы процентов) & ORC/Parquet иногда быстрее (единицы процентов) & чистая агрегация в JVM, для Vortex --- через границу JNI; нет проталкивания агрегатов в интерфейсе Paimon \\ \hline
    Хвосты задержки под конкурентной записью (среднее против медианы) & Vortex существенно у́же (среднее $\approx$ медиана, у Parquet среднее на $\approx 21\,\%$ выше медианы) & оба формата дают узкие хвосты разными способами & пул соединений Hadoop S3A эпизодически насыщается на сбросе чекпойнта; Vortex изолирован независимым нативным пулом, ORC сокращает число обращений к хранилищу за счёт крупных stripes и пропуска по индексным группам \\ \hline
\end{xltabular}
}

Если оставить только устойчивые наблюдения и отсечь запросы, на которых разница
не превышает разброса, преимущество Vortex над ORC сосредотачивается в трёх
классах запросов: запросы с селективным предикатом (проталкивание в нативный
декодер до границы JNI), диапазонные запросы по строкам с регулярной структурой
(FSST + пропуск по полным статистикам) и многошаговая аналитика с несколькими
предикатами и проходами. На простой группировке по столбцу низкой кардинальности
и на многомерной агрегации без предиката ORC оказывается не хуже Vortex или
лучше; на точечных запросах форматы сопоставимы.

Помимо распределения побед по типам запросов, обращает на себя внимание разница
в поведении форматов между фазами FINAL и DURING. По медиане ORC и Parquet
практически не меняют времени (отношение DURING/FINAL близко к единице), а
Vortex в DURING медленнее самого себя в FINAL на $\approx 22\,\%$. Это прямое
следствие зависимости специализированных кодеков Vortex (прежде всего FSST) от
порядка строк внутри блока: в свежих файлах уровня~0, формируемых по чекпойнту,
данные ещё не отсортированы по ключу слиянием, и плотность кодирования снижается;
после остановки записи и работы фонового слияния файлы становятся отсортированными
по ключу, и кодеки работают в полную плотность. ORC и Parquet таким эффектом не
обладают --- их кодеки менее чувствительны к порядку, --- поэтому их типичная
задержка между фазами и не меняется.

Сопоставление среднего и медианы по DURING-фазе позволяет оценить «хвосты»
задержки. У ORC и Vortex среднее и медиана практически совпадают (узкие хвосты),
у Parquet среднее на $\approx 21\,\%$ выше медианы из-за нескольких крупных
выбросов в десятки секунд в моменты сброса данных по чекпойнту. Архитектурное
объяснение --- общий пул соединений слоя совместимости Hadoop S3A эпизодически
насыщается под нагрузкой; Vortex от него изолирован независимым нативным пулом,
ORC компенсирует крупными полосами (stripes) и пропуском по индексным группам
(подробный разбор приведён в подразделе с результатами near-real-time-сценария).

Сопоставление двух серий экспериментов позволяет объяснить, почему Vortex
доминирует на TPC-DS, но не доминирует на потоковом LSM-сценарии. На TPC-DS
таблицы только для добавления, файлы крупные, путь чтения прямой колоночный,
накладные расходы границы JNI амортизируются на больших пакетах, а запросы ---
с высокой селективностью. На потоковом сценарии таблица --- с первичным ключом
и слиянием при чтении: итератор слияния открывает много файлов (особенно свежих,
не слитых файлов уровня~0), накладные расходы границы JNI накапливаются на
пер-файловой основе, данные в свежих файлах не упорядочены по ключу --- кодеки
Vortex (прежде всего FSST) теряют часть эффективности, --- а селективность
запросов в среднем ниже. Поэтому преимущество Vortex здесь скромнее и проявляется
в первую очередь на упорядоченном итоговом снимке и на запросах с селективным
предикатом. Областью наибольшей применимости Vortex остаётся именно batch-аналитика
на таблицах только для добавления, а не LSM с первичным ключом.

Преимущество Parquet на агрегации при полном сканировании объясняется отсутствием
проталкивания агрегатов в интерфейсе файловых форматов Apache Paimon в принципе
--- это не специфика Vortex. Агрегаты вычисляются движком Flink над
декодированными данными. Для Parquet весь конвейер --- декодер и агрегатор ---
находится в JVM, границы нет. Для Vortex данные декодируются нативно, затем
пересекают границу JNI (без копирования, но с фиксированной стоимостью на пакет)
и агрегируются в JVM; на миллионах пакетов это даёт заметную добавку. Устранение
ограничения --- проталкивание агрегатов в нативный SIMD-слой Vortex --- естественное
направление развития Apache Paimon и потенциальный вклад в его экосистему.

\subsection{Рекомендации по применению форматов хранения}

Рекомендации сформулированы с учётом сценария использования.

1. Batch-аналитика, таблицы только для добавления. Это основная
область применимости Vortex. На наборе TPC-DS масштаба \SI{100}{\giga\byte} Vortex
--- быстрейший из четырёх форматов (быстрее всех на \num{90} запросах из \num{104}),
по суммарному времени опережает Parquet на \SI{25}{\percent}, ORC --- на
\SI{32}{\percent}, пиковое ускорение --- до \SI{83}{\percent} (запросы с селективным
предикатом, узкой проекцией, многократным сканированием немногих столбцов). Условия
этого сценария --- крупные файлы, отсутствие слияния версий, упорядоченные данные,
высокая селективность --- максимально благоприятны для специализированных кодеков и
нативного проталкивания предиката. Рекомендуется выбирать Vortex как
формат хранения данных для аналитических витрин на основе таблиц только для
добавления (фактовые таблицы, журналы, исторические данные, материализованные
выборки).

2. Near-real-time-витрина (таблица с первичным ключом, потоковая запись).
По типичной задержке (медиана) Vortex --- быстрейший формат в обеих фазах
и на большинстве запросов (девять из тринадцати); рекомендуется, если в нагрузке
преобладают запросы с селективным предикатом (отсечение до агрегации), диапазонные
запросы по строкам с регулярной структурой (идентификаторы, пути, URL) и
многошаговая аналитика (последовательности событий, соединения с диапазонными
фильтрами). ORC --- разумный выбор, если в нагрузке преобладают простая
группировка по столбцу низкой кардинальности, чистая многомерная агрегация без
предиката или если важна узкая задержка отдельных сэмплов (наихудший случай, p99) ---
ORC демонстрирует одинаково узкие хвосты с Vortex, а в отдельных запросах ему
способствует словарное кодирование. ORC также даёт самый компактный объём хранения
(\SI{17}{\giga\byte} против \SI{18}{\giga\byte} у Parquet и \SI{20}{\giga\byte} у
Vortex на потоковом сценарии; на TPC-DS преимущество ORC меньше). Parquet
--- разумный универсальный выбор по умолчанию: не уступает по записи, хуже по
медиане чтения, зрелый и широко поддержанный; основной недостаток --- эпизодические
выбросы времени запросов под конкурентной записью в моменты сброса данных по
чекпойнту (среднее заметно выше медианы).

3. Что не рекомендуется. Lance в текущем состоянии интеграции с
Apache Paimon не рекомендуется для потоковых таблиц с первичным ключом (требуется
доработка слияния версий); на пакетной аналитике TPC-DS он оказался самым медленным
из четырёх форматов (в среднем в 1{,}5--2 раза медленнее по времени запросов и в
8 раз больше по объёму хранилища). Область Lance --- нагрузки машинного обучения с
произвольным доступом к строкам и векторным поиском, выходящие за рамки настоящей
работы. Apache Avro ---
строчный формат, в принципе непригоден для OLAP-нагрузки; применение его для
свежих файлов уровня~0 в LSM-дереве было проверено и отклонено --- аналитические
запросы Vortex на итоговом снимке замедляются за счёт совмещения построчных дельт
с колоночными файлами в итераторе слияния.

4. Конфигурация (общие соображения для NRT-витрин). Распределять слоты
между записью и чтением так, чтобы читатель не блокировался; увеличить размер
пакета записи Vortex до десятков тысяч строк (снижает накладные расходы JNI при
чтении в десятки раз); включить полные статистики по столбцам
(\texttt{metadata.stats-mode = full}); настроить расширенные политики хранения
снимков и манифестов под частые фиксации; не применять агрессивных тонких
настроек размера файлов и страниц (благоприятствуют Parquet и раздувают
хранилище) и не использовать строчный формат для свежих файлов уровня~0;
увеличивать интервал чекпойнтов следует с учётом того, что это укрупняет всплеск
ввода-вывода при сбросе.
